% Fallbacks and common problems of the Day 4 lab. The rows come from
% Labs/day04/README.md.
\begin{frame}{If a part does not work, and common problems}
  \footnotesize
  \renewcommand{\arraystretch}{1.2}
  \setlength{\tabcolsep}{4pt}
  \begingroup\centering
  \begin{tabular}{@{}L{0.40\textwidth}L{0.55\textwidth}@{}}
    \toprule
    \textbf{Problem} & \textbf{Fix or fallback} \\
    \midrule
    Your group has no dataset of Day 2
      & Copy the folder \code{data/} of a different group into
        \code{Labs/day02/}. Write this in the report. \\
    Task 4 or task 5 is not complete in time
      & Copy the two functions from the notebook of \code{solutions/}.
        Parts C and D need them. \\
    The \code{int8} model has an accuracy near 25 percent
      & The test code gives zeros or float values to the model. Check
        task 5. \\
    \code{Task 7: not complete}, the gradient is \code{None}
      & The function returns \code{x\_q}. Return
        \code{x + ops.stop\_gradient(x\_q - x)}. \\
    The build prints \code{model\_int8.h: No such file}
      & Run section 15 of the notebook, or use the files of the
        repository. \\
    The board gives 20 correct windows of 80
      & The output is 0 for each class. Complete task D2. \\
    \code{same class as on the laptop} is below 80
      & Check the scale and the zero point in tasks D1 and D2. Run section
        15 again. \\
    Tasks D1 and D2 are not complete in time
      & Use the sketch of \code{solutions/} with your four files of the
        notebook. \\
    \bottomrule
  \end{tabular}\par\endgroup
  \vskip 0.4em
  \small
  The README of the lab has the complete list.
\end{frame}
